
\documentclass[aps,preprint]{revtex4}
%%%%%%%%%%%%%%%%%%%%%%%%%%%%%%%%%%%%%%%%%%%%%%%%%%%%%%%%%%%%%%%%%%%%%%%%%%%%%%%%%%%%%%%%%%%%%%%%%%%%%%%%%%%%%%%%%%%%%%%%%%%%%%%%%%%%%%%%%%%%%%%%%%%%%%%%%%%%%%%%%%%%%%%%%%%%%%%%%%%%%%%%%%%%%%%%%%%%%%%%%%%%%%%%%%%%%%%%%%%%%%%%%%%%%%%%%%%%%%%%%%%%%%%%%%%%
\usepackage{amsfonts}
\usepackage{amsmath}
\usepackage{amssymb}
\usepackage{graphicx}

\setcounter{MaxMatrixCols}{10}
%TCIDATA{OutputFilter=LATEX.DLL}
%TCIDATA{Version=5.50.0.2960}
%TCIDATA{<META NAME="SaveForMode" CONTENT="1">}
%TCIDATA{BibliographyScheme=Manual}
%TCIDATA{Created=Friday, April 06, 2012 11:15:25}
%TCIDATA{LastRevised=Friday, January 16, 2026 09:09:42}
%TCIDATA{<META NAME="GraphicsSave" CONTENT="32">}
%TCIDATA{<META NAME="DocumentShell" CONTENT="Articles\SW\REVTeX 4">}
%TCIDATA{Language=American English}
%TCIDATA{CSTFile=revtex4.cst}

\newtheorem{theorem}{Theorem}
\newtheorem{acknowledgement}[theorem]{Acknowledgement}
\newtheorem{algorithm}[theorem]{Algorithm}
\newtheorem{axiom}[theorem]{Axiom}
\newtheorem{claim}[theorem]{Claim}
\newtheorem{conclusion}[theorem]{Conclusion}
\newtheorem{condition}[theorem]{Condition}
\newtheorem{conjecture}[theorem]{Conjecture}
\newtheorem{corollary}[theorem]{Corollary}
\newtheorem{criterion}[theorem]{Criterion}
\newtheorem{definition}[theorem]{Definition}
\newtheorem{example}[theorem]{Example}
\newtheorem{exercise}[theorem]{Exercise}
\newtheorem{lemma}[theorem]{Lemma}
\newtheorem{notation}[theorem]{Notation}
\newtheorem{problem}[theorem]{Problem}
\newtheorem{proposition}[theorem]{Proposition}
\newtheorem{remark}[theorem]{Remark}
\newtheorem{solution}[theorem]{Solution}
\newtheorem{summary}[theorem]{Summary}
\newenvironment{proof}[1][Proof]{\noindent\textbf{#1.} }{\ \rule{0.5em}{0.5em}}
\input{tcilatex}
\begin{document}

\preprint{}
\title{Complex Symmetric Matrices}
\author{Brian Weiner}
\affiliation{PSU}
\author{}

\begin{abstract}
\end{abstract}

\date{07/19/2025}
\startpage{1}
\maketitle
\tableofcontents

\section{Theorem on Complex Antisymmetric Matrices}

\begin{theorem}
For any \textbf{complex antisymmetric matrix }$A=B+iC$ where:

\begin{itemize}
\item $B^{T}=-B$ (real part of antisymmetric matrix) $=-\func{Re}\left(
A\right) $

\item $C^{T}=-C$ (complex part of antisymmetric matrix) $=-\func{Im}\left(
A\right) $
\end{itemize}

the following conditions hold:

\begin{enumerate}
\item $A$ is normal ($AA^\dagger = A^\dagger A$) if and only if the
commutator $[B,C] = BC - CB = 0$

\item When normal, $A$ is unitarily diagonalizable with purely imaginary
eigenvalues

\item When $[B,C] \neq 0$, $A$ may have defective eigenvalues requiring
Jordan form
\end{enumerate}
\end{theorem}

\begin{proof}
\underline{Part 1: Normality Condition}

For $A = B + iC$ with $B^T = -B$ and $C^T = -C$:

\begin{align*}
AA^\dagger &= (B + iC)(B^T - iC^T) \\
&= (B + iC)(-B - iC) \\
&= -B^2 - iBC - iCB + C^2
\end{align*}

\begin{align*}
A^\dagger A &= (B^T - iC^T)(B + iC) \\
&= (-B - iC)(B + iC) \\
&= -B^2 - iCB - iBC + C^2
\end{align*}

The normality condition reduces to: 
\begin{equation}
AA^\dagger - A^\dagger A = -2i[B,C] = 0 \quad \iff \quad [B,C] = 0
\end{equation}

\underline{Part 2: Diagonalizability}

When $[B,C] = 0$:

\begin{itemize}
\item $A$ is normal $\Rightarrow$ unitarily diagonalizable via $A =
UDU^\dagger$

\item Eigenvalues are purely imaginary since $A^\dagger = -A$
\end{itemize}

When $[B,C] \neq 0$:

\begin{itemize}
\item $A$ is non-normal

\item May have defective eigenvalues (geometric multiplicity $<$ algebraic
multiplicity)

\item Requires Jordan form $A = PJP^{-1}$ with non-trivial Jordan blocks
\end{itemize}
\end{proof}

\section{Main Theorem}

\begin{theorem}
For an $n \times n$ complex matrix $A$, an eigenvalue $\lambda$ is \textbf{%
defective} if and only if its geometric multiplicity is strictly less than
its algebraic multiplicity: 
\begin{equation*}
\dim \ker(A - \lambda I) < \text{mult}_{\text{alg}}(\lambda).
\end{equation*}
\end{theorem}

\section{Proof}

The Jordan decomposition $A = PJP^{-1}$ yields:

\begin{equation}
J=\bigoplus_{i=1}^{k}J_{m_{i}}(\lambda _{i}),\quad J_{m}(\lambda )=%
\begin{pmatrix}
\lambda  & 1 &  & 0 \\ 
0 & \ddots  & \ddots  &  \\ 
& \ddots  & \ddots  & 1 \\ 
0 &  & 0 & \lambda 
\end{pmatrix}%
_{m\times m}
\end{equation}

For each $\lambda$:

\begin{itemize}
\item Algebraic multiplicity $\left( AM\right) $: $\sum_{m_{i}}m_{i}$ where $%
\lambda _{i}=\lambda $

\item Geometric multiplicity $\left( GM\right) $: number of blocks $%
J_{m_{i}}(\lambda )$
\end{itemize}

Defectiveness occurs when any $m_i > 1$, creating non-trivial Jordan blocks.

\section{Examples}

\subsection{Non-Defective Matrix}

\begin{equation}
A = 
\begin{pmatrix}
2 & 0 \\ 
0 & 2%
\end{pmatrix}%
\end{equation}

\begin{itemize}
\item $\chi _{A}(\lambda )=(\lambda -2)^{2}$ (AM $=2$)

\item $\ker (A-2I)=\mathbb{C}^{2}$ (GM $=2$)

\item Full set of eigenvectors
\end{itemize}

\subsection{Defective Matrix}

\begin{equation}
B = 
\begin{pmatrix}
2 & 1 \\ 
0 & 2%
\end{pmatrix}%
\end{equation}

\begin{itemize}
\item $\chi _{B}(\lambda )=(\lambda -2)^{2}$ (AM $=2$)

\item $\ker (B-2I)=\text{span}\{(1,0)^{T}\}$ (GM$\ =1$)

\item Requires Jordan form
\end{itemize}

\section{Antisymmetric Case}

For $A = B + iC$ with $B^T = -B$, $C^T = -C$:

\begin{equation}
\text{Normality} \Leftrightarrow [B,C] = 0
\end{equation}

Non-normal cases may exhibit defective eigenvalues:

\begin{equation}
A = 
\begin{pmatrix}
0 & 1+i \\ 
-1-i & 0%
\end{pmatrix}%
, \quad [B,C] \neq 0
\end{equation}

\begin{itemize}
\item Potential defective eigenvalues when $[B,C] \neq 0$

\item Block-diagonalization may contain non-trivial Jordan blocks
\end{itemize}

\section*{Conclusion}

Defective eigenvalues emerge when geometric multiplicity fails to match
algebraic multiplicity, necessitating Jordan forms. For complex
antisymmetric matrices, non-commutativity of real and imaginary components
can induce such defects, preventing full diagonalization.

\section{Block Diagonalization Theorem}

\begin{theorem}
Every \textbf{complex} antisymmetric matrix $A\in M_{n}(\mathbb{C})$ with $%
A^{T}=-A$ can be brought to block diagonal form: 
\begin{equation}
A=U%
\begin{pmatrix}
\begin{pmatrix}
0 & \lambda _{1} \\ 
-\lambda _{1} & 0%
\end{pmatrix}
&  &  &  &  &  \\ 
& \ddots  &  &  &  &  \\ 
&  & 
\begin{pmatrix}
0 & \lambda _{k} \\ 
-\lambda _{k} & 0%
\end{pmatrix}
&  &  &  \\ 
&  &  & 0 &  &  \\ 
&  &  &  & \ddots  &  \\ 
&  &  &  &  & 0%
\end{pmatrix}%
U^{T}
\end{equation}%
where:

\begin{itemize}
\item $U$ is unitary ($U^\dagger U = I$)

\item $\lambda _{j}\in \mathbb{C}$ are non-zero \textbf{complex} numbers

\item The last block is a $1 \times 1$ zero if $n$ is odd
\end{itemize}
\end{theorem}

\begin{proof}
We proceed through four main steps:

\textbf{Step 1}: Reduction to Singular Value Form

For any complex square matrix $A$, there exist unitary matrices $U$ and $V$
such that: 
\begin{equation}
A = U \Sigma V^T
\end{equation}
where $\Sigma$ is diagonal with non-negative real entries. For antisymmetric 
$A$: 
\begin{equation}
A^T = V \Sigma^T U^T = -A = -U \Sigma V^T
\end{equation}
This implies $V \Sigma U^T = -U \Sigma V^T$, leading to: 
\begin{equation}
\Sigma = -Q \Sigma Q^T \quad \text{where} \quad Q = U^\dagger V
\end{equation}

\textbf{Step 2}: Pairing of Singular Values

From $\Sigma = -Q \Sigma Q^T$, we deduce that non-zero singular values must
come in pairs. For each $\sigma_j > 0$, there exists $\sigma_k = \sigma_j$
such that the corresponding blocks satisfy: 
\begin{equation}
\begin{pmatrix}
\sigma_j & 0 \\ 
0 & \sigma_k%
\end{pmatrix}
= -%
\begin{pmatrix}
q_{11} & q_{12} \\ 
q_{21} & q_{22}%
\end{pmatrix}
\begin{pmatrix}
\sigma_j & 0 \\ 
0 & \sigma_k%
\end{pmatrix}
\begin{pmatrix}
q_{11} & q_{21} \\ 
q_{12} & q_{22}%
\end{pmatrix}%
\end{equation}
This enforces the $2 \times 2$ block structure.

\textbf{Step 3}: Construction of Block Diagonal Form

Through careful choice of unitary transformations, we can arrange: 
\begin{equation}
\Sigma = \bigoplus_{j=1}^k 
\begin{pmatrix}
0 & \sigma_j \\ 
-\sigma_j & 0%
\end{pmatrix}
\oplus 0_{n-2k}
\end{equation}
The complex phases of $\sigma_j$ can be absorbed into $U$, giving the
general form with $\lambda_j \in \mathbb{C}$.

\textbf{Step 4}: Unitary Congruence

The final form is achieved through unitary congruence (not similarity): 
\begin{equation}
A = U \Lambda U^T
\end{equation}
where $\Lambda$ contains the $2 \times 2$ blocks and possibly a zero block.
This preserves the antisymmetry property while revealing the canonical
structure.
\end{proof}

\section*{Corollaries}

\begin{corollary}
For \textbf{normal} antisymmetric matrices ($[B,C]=0$), all $\lambda _{j}$
are purely imaginary.
\end{corollary}

\begin{corollary}
When $n$ is odd, at least one eigenvalue must be zero.
\end{corollary}

\end{document}
